\subsection{章末確認問題}
\begin{itemize}[leftmargin=2em]
\item 命題と述語の違いを、例を挙げて説明せよ。
\item 「すべての利用者は利用者IDを持つ」を述語論理の考え方で説明せよ。
\item 背理法と対偶証明は、どのような場面で使いやすいか。
\item 集合の共通部分、和集合、差集合を、権限管理の例で説明せよ。
\item 関係と関数の違いを説明せよ。API設計でこの違いが重要になる例を述べよ。
\item 有向グラフと重み付きグラフが、ソフトウェア開発で使われる例を一つずつ挙げよ。
\item 二分探索木を中順に巡回すると、なぜ昇順の列が得られるかを説明せよ。
\item 有限状態機械が画面遷移や通信プロトコルの設計に向く理由を説明せよ。
\item 正規文法、文脈自由文法、正規表現の関係を説明せよ。
\item 剰余演算がソフトウェアで使われる例を三つ挙げよ。
\item 全組合せテストが現実的でない場合、数え上げの知識がどう役立つかを説明せよ。
\item 離散確率分布において、各確率の合計が1でなければならない理由を説明せよ。
\item 正確度と精密度の違いを説明せよ。
\item モノイドが分散処理や集計に役立つ理由を説明せよ。
\item 微分と積分が機械学習やシミュレーションで使われる理由を説明せよ。
\end{itemize}
\begin{notebox}{解答の方向性}
1は真偽が決まる文と、変数を含む性質・関係の違いである。4は集合演算を利用者集合や権限集合に対応させる。8は「状態」と「入力」と「遷移」で振る舞いを表にできる点を述べる。13は、真値への近さと測定値同士の近さの違いを述べる。14は、結合演算と単位元があるため部分結果を安全に合成しやすい点を述べる。
\end{notebox}
